\section{Discussion}

\subsection{The Performance Variation Hierarchy}

The fixed-configuration analysis yielded a clear hierarchy:

\begin{center}
Data source heterogeneity $\gg$ Splitting method $\approx$ In-domain algorithm choice
\end{center}

The source heterogeneity gap ($\Delta = 0.390$) substantially exceeded the 
split-method gap ($\Delta = 0.046$). The split-method gap, in turn, was comparable 
to the in-domain algorithm spread (MCC range = 0.063). Three caveats apply. First, the 
hierarchy was established on a single Ames benchmark. Second, the LDO $\Delta$ reflects 
differences in test set size and class distribution across directions and should therefore 
be treated as an upper bound. Third, algorithm robustness varied under source shift; 
algorithm choice may thus matter more in cross-source deployment than the in-domain 
figures suggest.

The LDO experiment should be interpreted as a diagnostic stress test, not a realistic 
deployment scenario. In practice, models are trained on all available data regardless of 
provenance. Mutagenic potential is a property of molecular structure, not database origin. 
A large LDO drop suggests that apparent in-domain performance depends strongly on the 
joint composition of the training and test sets. This pattern suggests incomplete learning 
of broadly generalizable structure-activity relationships. 
This is a concern for regulatory applications: new chemicals registered under EU REACH or 
K-REACH may occupy regions of chemical space that are underrepresented in pharmaceutical 
benchmarks, and a benchmark MCC of 0.67 on a mixed-source test set provides limited 
assurance in that context. Cross-source LDO reveals this dependence and offers a more 
conservative estimate of out-of-distribution performance.

\subsection{What Drives the Data Source Heterogeneity Gap?}

Four non-exclusive factors may contribute, listed from most to least parsimonious:

\begin{enumerate}
    \item \textbf{Prior shift.} The 19.5-fold difference in positive rates between sources 
    miscalibrates the default decision threshold. This is consistent with XGB achieving a 
    sensitivity of only 0.062 in the industrial$\to$drug direction.

    \item \textbf{Data quality and labeling differences.} Drug-sourced compounds come from 
    curated benchmarks; industrial compounds come from regulatory submissions (ECHA, 
    K-REACH) with varying protocols. Prevalence differences may therefore arise independently 
    of underlying chemical differences.

    \item \textbf{Selection bias.} Drug candidates are optimized for biological activity, 
    while industrial chemicals are selected for functional or commercial utility. These 
    represent fundamentally different sampling criteria.

    \item \textbf{Chemical space divergence.\cite{Lipinski_2004,Veber_2002}} Even a 
    perfectly calibrated model may generalize poorly if the two source populations occupy 
    distinct regions of chemical space.
\end{enumerate}

To estimate the contribution of prior shift, each LDO direction was evaluated at three 
thresholds: \textit{default} ($t = 0.5$), \textit{prior-corrected} ($t$ set to the 
training positive rate, optimal under pure prior shift), and \textit{oracle} ($t$ 
maximizing MCC on the test set, an upper bound on recalibration). This analysis used 
Morgan FP-1024 plus 13 physicochemical descriptors rather than the full 138-feature compact 
set, so absolute MCC values differ from Table~\ref{tab:ldo} (notably RF 
industrial$\to$drug: 0.102 here vs.\ 0.325 in the main analysis). The three-threshold 
comparison within each row is nonetheless internally valid.

The results (Table~S7) show that prior-shift correction gave modest gains (mean $+0.075$; 
range $-0.024$ to $+0.239$). Oracle MCC topped out at 0.27--0.36, well below the in-domain 
scaffold-CV MCC of 0.624. This leaves a residual gap of 0.26--0.35 that threshold 
adjustment alone cannot close. Sensitivity to threshold choice also varied by algorithm: 
RF gained $+0.243$ MCC from default to oracle in the industrial$\to$drug direction, while 
gradient boosting models gained only $+0.07$ to $+0.10$. Prior shift thus accounts for 
part of the LDO performance loss, but a substantial residual gap remains. The practical 
implication holds regardless of the underlying cause: published Ames QSAR performance is highly 
sensitive to source composition, and recalibration alone is insufficient under cross-source 
evaluation.

\subsection{\textit{In Vivo} Micronucleus and Tiered Testing}

\textit{In vivo} MN models showed strong ranking performance (ROC-AUC = 0.860, PR-AUC = 0.237), 
but classification at any fixed threshold was unreliable (MCC 95\% CI crossed zero). This 
indicates that current 2D-descriptor models are not suitable for regulatory binary hazard 
decisions under ICH M7\cite{ICH_M7_2017} or REACH.\cite{EU_REACH_2006} The 8.5-fold 
enrichment over baseline prevalence, however, suggests potential utility as a prioritization 
tool within a tiered screening workflow.\cite{Dobo_2012}

Performance was also highly sensitive to preprocessing: salt stripping improved MCC by 
$+0.22$. These findings indicate that reported \textit{in vivo} performance can differ appreciably depending 
on undocumented preprocessing choices.\cite{Fourches_2016} Mechanistically, micronucleus 
induction is governed by ADME, metabolic activation, tissue distribution, and DNA 
repair\cite{Tweats_2007,Kirkland_2011}---processes that 2D descriptors do not represent.

\subsection{Cross-Endpoint Overlap and Preprocessing}

More than 90\% of the \textit{in vitro} and \textit{in vivo} compounds also appear in the Ames dataset. 
Performance differences across endpoints (Ames MCC = 0.555 vs.\ \textit{in vivo} 0.311) therefore 
most likely reflect intrinsic endpoint predictability from 2D structure, not differences in 
chemical coverage.\cite{Tropsha_2010} Multi-task learning may be a productive avenue for 
future work.\cite{Ramsundar_2015}

Preprocessing effects were endpoint-specific: salt stripping improved \textit{in vivo} performance, 
produced mixed results for Ames, and reduced \textit{in vitro} performance. The reduction for 
\textit{in vitro} may reflect clastogenic activity of some counterions. Although standardized 
preprocessing pipelines are widely recommended,\cite{Fourches_2010,Fourches_2016} the 
present results indicate that endpoint-specific optimization is preferable.

\subsection{Practical Recommendations}

The following reporting standards are proposed for genotoxicity QSAR:

\begin{enumerate}
    \item Use scaffold-based splitting.
    \item When source labels are available, report cross-source LDO performance with 
    disclosed source composition.
    \item Optimize preprocessing separately for each endpoint.
    \item Report bootstrap CIs and applicability domain coverage.
    \item For class-imbalanced endpoints, report both threshold-free (ROC-AUC, PR-AUC) and 
    threshold-dependent (MCC, balanced accuracy) metrics.
    \item Test algorithm superiority claims under source shift, not only in-domain.
\end{enumerate}

Several limitations apply. Only a single scaffold split was used. The GCN lacked edge 
features and pretraining.\cite{Hu_2020_pretrain} Three-dimensional descriptors, metabolite 
predictions, within-AD vs.\ out-of-AD stratification for LDO, and multiple-comparison 
correction for McNemar tests\cite{Demsar_2006} were not evaluated.

